\section{Metrics and Conventions}
\label{sec:metrics}

\subsection{Fault quantities}

All four metrics are built from the short-circuit power and the impedance matrix at the evaluation
bus. For bus $i$ in a given network and operating state, the short-circuit power $\mathrm{SCMVA}_i$ (MVA) is
obtained from a bolted three-phase fault solved from the actual pre-fault power flow, with
$\mathrm{SCMVA}_i = S_\mathrm{b}|I_{\mathrm{f},i}|$ at $V_\mathrm{n} = 1$~pu and $S_\mathrm{b} = 100$~MVA. The
impedance matrix $\mathbf{Z} = \mathbf{Y}^{-1}$ comes from the same fault model: the passive
network, the source admittances of the committed synchronous generators, and constant-impedance
equivalents of the loads at their pre-fault voltages. Assuming constant-impedance loads and constant-current
converters during the fault, the fault network is linear, and
$\mathrm{SCMVA}_i \approx S_\mathrm{b}|V_i|/|Z_{ii}|$: the short-circuit power carries both the driving-point
impedance $Z_{ii}$ and the pre-fault voltage magnitude $|V_i|$. Every denominator uses the IBR plants' rated active
power $P_i$.

\subsection{Metrics}

Table~\ref{tab:conventions} lists the four metrics, each with the network, evaluation point and
interaction term its source specifies; Section~\ref{sec:study} gives the plant model and where each
point lies in it. SCR and WSCR are evaluated at the point of interconnection (POI), on the
transmission network without the IBR plants \cite{ieee2800}; WSCR combines the three sites into one
group value \cite{ieee2800,zhang2014}. ESCR is evaluated at each plant's medium-voltage (MV) bus,
the high-voltage side of its padmount transformer, on the network with the IBR plants
\cite{fingrid2023}. It is the multi-infeed ESCR of
\cite{cigre364} with no reactive compensation, the quantity titled SCRIF in \cite{ieee2800} and
equivalent SCR in \cite{cigre671,fingrid2023}; it is not the effective SCR of line-commutated HVDC,
which subtracts filter reactive power. SDSCR is evaluated at the converter terminal, on the low-voltage
side of the padmount transformer, as in \cite{wu2018}. Its weights $w_{ij}$ are complex, with the conjugated
voltage ratio as in \cite{wu2018,maharjan2023}, and
the magnitude is taken of the whole weighted sum.

\begin{table}[!t]
\caption{Metric definitions and conventions}
\label{tab:conventions}
\centering
\footnotesize
\setlength{\tabcolsep}{3pt}
\begin{tabular}{@{}llllc@{}}
\toprule
Metric & Definition & IBR & Point & Ref. \\
\midrule
SCR & $S_i / P_i$ & no & POI & \cite{ieee2800} \\
WSCR & $\sum_i S_i P_i \,/\, (\sum_i P_i)^2$ & no & POI & \cite{ieee2800} \\
ESCR & $S_i \,/\, (P_i + \sum_{j \ne i} \mathrm{MIIF}_{ji} P_j)$ & yes & MV bus & \cite{cigre364} \\
SDSCR & $|V_i|^2 \,/\, (|Z_{ii}|\,|P_i + \sum_{j \ne i} w_{ij} P_j|)$ & yes & converter terminal & \cite{wu2018} \\
\midrule
\multicolumn{5}{@{}p{\linewidth}@{}}{\scriptsize $\mathrm{MIIF}_{ji} = |Z_{ji}/Z_{ii}|$; $w_{ij} = (Z_{ij}/Z_{ii})\,(V_i/V_j)^*$; $Q_f = 0$ in ESCR; $P$: rated active power.} \\
\bottomrule
\end{tabular}

\end{table}

\subsection{Interaction factor}

ESCR and SDSCR add an interaction term to the plant's own power: $\sum_{j\ne i}
\mathrm{MIIF}_{ji}P_j$ and $\sum_{j\ne i} w_{ij}P_j$, respectively. SCR and WSCR have none. The
multi-infeed interaction factor $\mathrm{MIIF}_{ji}$ is the ratio of the voltage changes
$\Delta V_j/\Delta V_i$ for a small voltage step at bus $i$ \cite{cigre364}. The step is about
1\,\%, small enough that no voltage-controlling device acts, and the ratio then equals the
impedance-matrix form $|Z_{ji}/Z_{ii}|$ \cite{cigre364}. This impedance form is the reference used
here. Other sources obtain the factor differently: IEEE Std 2800-2022 \cite{ieee2800} allows impedance-matrix
manipulation or reactive-power injection, and Fingrid \cite{fingrid2023} applies a voltage step in a power
flow, in which voltage-regulating generators respond. Section~\ref{sec:study} describes the
voltage-step variant used to test this choice.

\subsection{Attribution of changes}

All four metrics share one structure. With the approximation of Section~\ref{sec:metrics}-A,
SCR, ESCR and SDSCR can be written as
\begin{equation}
m \approx \frac{S_\mathrm{b}\,|V_i|^{a}}{|Z_{ii}|\,D_i} ,
\label{eq:shape}
\end{equation}
exactly for SDSCR, with the exponent $a$ and the denominator $D_i$ in Table~\ref{tab:shape}.
\begin{table}[!t]
\caption{Exponent and denominator of each metric in \eqref{eq:shape}}
\label{tab:shape}
\centering
\footnotesize
\renewcommand{\arraystretch}{1.3}
\begin{tabular}{@{}lcl@{}}
\toprule
Metric & $a$ & $D_i$ \\
\midrule
SCR & 1 & $P_i$ \\
ESCR & 1 & $P_i + \sum_{j\ne i}\mathrm{MIIF}_{ji}P_j$ \\
SDSCR & 2 & $|P_i + \sum_{j\ne i} w_{ij}P_j|$ \\
\bottomrule
\end{tabular}
\end{table}
WSCR is a weighted sum of the three sites' SCRs, with fixed weights $(P_i/\sum_k P_k)^2$. With $P_i$ the rated power in every state, the
change of a metric between two states (two hours, or an outage and the intact case) splits into
three parts:
\begin{equation}
\Delta\ln m = \underbrace{a\,\Delta\ln|V_i|}_{\delta_V}
\;\underbrace{-\,\Delta\ln D_i}_{\delta_I}
\;\underbrace{-\,\Delta\ln|Z_{ii}|}_{\delta_Z} ,
\label{eq:split}
\end{equation}
the pre-fault-voltage, interaction and driving-point-impedance parts,
$\delta_V = a\,\Delta\ln|V_i|$, $\delta_I = -\Delta\ln D_i$ and $\delta_Z = -\Delta\ln|Z_{ii}|$,
each positive when it raises the metric. The interaction part is zero
for SCR and WSCR. Since~\eqref{eq:shape} is approximate, the parts are computed exactly by recomputing the metric
with unit pre-fault voltage and without the interaction factor. A site is called weaker when its
metric falls, and the network weaker when $|Z_{ii}|$ rises ($\delta_Z < 0$).
